% ch07_u12 —— 原书 p.346–p.355（第 7 章 7.13.18 末–7.14.3）
% 底本：A. Beilinson, V. Drinfeld, Quantization of Hitchin's Integrable System
% and Hecke Eigensheaves
% 承接：上单元（p.336–p.345）末句断在 "...are c-lattices such that"，
% 本单元正文自「$A(F)\subset U$、$A(U')=0$，则得到……」续起。

\SourcePage{346}

$A(F)\subset U$、$A(U')=0$，则得到 $A^{\sim}:U/U'\to U/U'$，并令 $\operatorname{tr}A:=\operatorname{tr}A^{\sim}$）。

泛函 $\operatorname{tr}:E_{tr}\to\mathbb{C}$ 对 $E^{Lie}$ 的伴随作用不变；它在 $[E_+,E_-]\subset E_{tr}$ 上取零值。

我们的扩张 $E^{\flat}$ 带典范分裂 $s_+:E_+\to E^{\flat}$、$s_-:E_-\to E^{\flat}$。即：对 $A\in E_+$，其提升 $s_+(A)$ 由下述性质刻画——$s_+(A)$ 零化 $C$ 中一切被来自 $\operatorname{Im}A\subset\mathfrak{g}$ 的 Clifford 算子零化的元素。同样，$s_-(A)$ 是 $A\in E_-$ 的唯一提升，它零化 $C$ 中一切被来自 $(\operatorname{Ker}A)^{\perp}\subset F^{*}$ 的 Clifford 算子零化的元素。截面 $s_{\pm}$ 与 $E$ 的伴随作用交换；且对 $A\in E_{tr}$ 有 $s_-(A)-s_+(A)=\operatorname{tr}A\in\mathbb{C}\subset E^{\flat}$。
% [Tutor Note: 原书 p.346 此句印作 "from Im $A\subset\mathfrak{g}$"；因 $A\in E_+$ 是 $F$ 的有界自同态，此处 $\mathfrak{g}$ 疑为 $F$ 之误，底本如此，保留原样。]

易见具有这些性质的资料 $(E^{\flat},s_{\pm})$ 唯一确定。确实，考虑 $E$-双模的精确序列
\begin{equation}\tag{373}
0\to E_{tr}\xrightarrow{(-,+)}E_+\oplus E_-\xrightarrow{(+,+)}E\to 0.
\end{equation}
现在 $s=(s_+,s_-)$ 把 $E^{\flat}$ 与扩张 (373) 沿 $\operatorname{tr}:E_{tr}\to\mathbb{C}$ 的前推等同。$E^{Lie}$ 在 $E^{\flat}$ 上的伴随作用来自 $E$-双模 $E_+\oplus E_-$ 上的伴随作用。

\emph{注.} (i) 向量空间 $F\otimes F^{*}$ 上有 4 个自然拓扑，其开子空间基分别由 $U\otimes V$、$U\otimes F^{*}$、$F\otimes V$ 与 $U\otimes F^{*}+F\otimes V$ 构成，其中 $U\subset F$、$V\subset F^{*}$ 是开子空间。相应的完备化分别等于 $E_{tr}$、$E_+$、$E_-$ 与 $E$。迹泛函是典范配对 $F\otimes F^{*}\to\mathbb{C}$ 的连续延拓。

(ii) 置 $(E_-/E_{tr})^{\flat}:=E_-/\operatorname{Ker}\operatorname{tr}$；这是 $(E_-/E_{tr})^{Lie}$ 被 $\mathbb{C}$ 的中心扩张。注意 $E_-/E_{tr}\cong E/E_+$，故有投影 $\pi_-:E^{Lie}\to(E_-/E_{tr})^{Lie}$。它典范地提升为扩张的态射 $\pi^{\flat}_-:E^{\flat}\to(E_-/E_{tr})^{\flat}$，其核为 $s_+(E_+)$。换言之，$E^{\flat}$ 是 $(E_-/E_{tr})^{\flat}$ 沿 $\pi_-$ 的拉回。$\pm$ 互换时同理。

(iii) 设 $F^{i}$ 是 $F$ 由闭子空间给出的有限滤链；以 $B\subset E_{F}$ 记保持该滤链的自同态构成的子代数。我们有\SourcePage{347}%
$B^{Lie}$ 的诱导中心扩张 $B^{\flat}$。另一方面，有显然的投影 $\operatorname{gr}^{i}:B\to E_{\operatorname{gr}^{i}F}$；以 $B^{\flat i}$ 记扩张 $E^{\flat}_{\operatorname{gr}^{i}F}$（即 $E^{Lie}_{\operatorname{gr}^{i}F}$ 的中心扩张）沿该投影的拉回。以 $B^{\flat'}$ 记诸扩张 $B^{\flat i}$ 的 Baer 和。则存在典范（且唯一）的扩张同构 $B^{\flat'}\xrightarrow{\sim}B^{\flat}$。

确实，$B^{\flat'}$ 与 (373) 的精确子序列
\[
0\to B\cap E_{tr}\to(B\cap E_+)\oplus(B\cap E_-)\to B\to 0
\]
所定义的扩张一致（注意对 $e\in B\cap E_{tr}$ 有 $\operatorname{tr}(e)=\sum\operatorname{tr}(\operatorname{gr}^{i}e)$）。特别地我们看到，$B^{\flat}$ 在李子代数 $\operatorname{Ker}\operatorname{gr}^{\cdot}$ 上典范分裂。

\bdp{7.13.19} 置 $K=\mathbb{C}((t))$、$O:=\mathbb{C}[[t]]$。设 $F$ 是带通常拓扑的有限维 $K$-向量空间；这是 Tate $\mathbb{C}$-向量空间。设 $i:\mathcal{D}\hookrightarrow E$ 是作用在 $F$ 上的 $K$-微分算子代数，于是得到 $\mathcal{D}^{Lie}$ 的中心扩张 $\mathcal{D}^{\flat}$。下面（沿用 [BS] 2.4）用几何的语言重述 $\mathcal{D}^{\flat}$ 的 Tate 描述。
% [Tutor Note: 原书 p.347 此处印作 "the agebra of K-differential operators"，"agebra" 系 "algebra" 之误，底本如此，保留原样。]

置 $F':=\operatorname{Hom}_{K}(F,K)$、$F^{\circ}:=F'\otimes_{K}\omega_{K}$。显然 $F^{\circ}$ 与 Tate 对偶 $F^{*}$ 一致（用配对 $f^{\circ},f\mapsto\langle f^{\circ},f\rangle:=\operatorname{Res}(f^{\circ},f)$）。我们的 $F$ 是左 $\mathcal{D}$-模，而 $F^{\circ}$ 带唯一的右 $\mathcal{D}$-模结构，使 $\langle,\rangle$ 为 $\mathcal{D}$-不变配对；注意 $\mathcal{D}$ 以微分算子作用在 $F^{\circ}$ 上，这就是通常几何中的「伴随」作用。

设 $K\widehat{\otimes}K$ 是 $K\otimes K$ 关于以 $(t^{n}O)\otimes(t^{n}O)$ 为基的拓扑之完备化，即 $K\widehat{\otimes}K:=\mathbb{C}[[t_{1},t_{2}]][t_{1}^{-1}][t_{2}^{-1}]$。设 $F\widehat{\otimes}F^{\circ}$ 是 $F\otimes F^{\circ}$ 的类似完备化；这是有限维 $K\widehat{\otimes}K$-模。以 $F\widehat{\otimes}F^{\circ}(\infty\Delta)$ 记 $F\widehat{\otimes}F^{\circ}$ 对 $(t_{1}-t_{2})^{-1}$（即对角线方程）的局部化。考虑标准精确序列
\begin{equation}\tag{374}
0\to F\widehat{\otimes}F^{\circ}\to F\widehat{\otimes}F^{\circ}(\infty\Delta)\xrightarrow{r}\mathcal{D}\to 0
\end{equation}
其中投影 $r$ 把「核」$k=k(t_{1},t_{2})dt_{2}\in F\widehat{\otimes}F^{\circ}(\infty\Delta)$ 送到微分算子 $r(k):F\to F$，$f(t)\mapsto\operatorname{Res}_{t_{2}=t}(k(t,t_{2}),f(t_{2}))dt_{2}$。注意 $F\widehat{\otimes}F^{\circ}$ 显然以自然的方式是 $\mathcal{D}$-双模。此双作用\SourcePage{348}%
唯一地延拓为 $F\widehat{\otimes}F^{\circ}(\infty\Delta)$ 上与 $K$-双模结构相容的 $\mathcal{D}$-双作用。易见 (374) 是 $\mathcal{D}$-双模的精确序列。设 $\operatorname{tr}:F^{\circ}\widehat{\otimes}F\to\mathbb{C}$ 为态射 $f\otimes f^{\circ}\mapsto\langle f^{\circ},f\rangle$（即它是对角线上限制的留数）。它关于 $\mathcal{D}^{Lie}$ 的伴随作用不变。以 $D^{\flat'}$ 记 (374) 沿 $\operatorname{tr}$ 的前推。$\mathcal{D}$ 在 $F\widehat{\otimes}F^{\circ}(\infty\Delta)$ 上的伴随作用给出 $D^{\flat'}$ 上的 $D^{Lie}$-模结构。对 $l^{\flat}_{1},l^{\flat}_{2}\in D^{\flat'}$，置 $[l^{\flat}_{1},l^{\flat}_{2}]:=l_{1}(l^{\flat}_{2})$，其中 $l_{1}$ 是 $l^{\flat}_{1}$ 在 $D^{Lie}$ 中的像。
% [Tutor Note: 原书 p.348 此处作 "The adjoint action of  on $F\hat\otimes F^{\circ}(\infty\Delta)$ yields a $D^{Lie}$-module structure on $D^{\flat'}$"，"of" 之后脱漏 $\mathcal{D}$；译文按上下文补出，底本如此，保留原样。]

\bdp{7.13.20} \emph{引理.} 括号 $[\,,\,]$ 斜对称，故它使 $D^{\flat'}$ 成为 $D^{Lie}$ 被 $\mathbb{C}$ 的中心扩张。存在中心扩张的唯一同构
\[
D^{\flat'}\xrightarrow{\sim}D^{\flat}.
\]

\emph{证明.} 只需建立 $D^{Lie}$-模扩张的同构 $D^{\flat'}\xrightarrow{\sim}D^{\flat}$。它来自 $\mathcal{D}$-双模精确序列的典范嵌入 $i^{\sim}:(374)\hookrightarrow(373)$，后者定义如下。态射 $\mathcal{D}\hookrightarrow E$ 是我们的标准嵌入 $i$，而 $i^{\sim}:F\widehat{\otimes}F^{\circ}=F\widehat{\otimes}F^{*}\xrightarrow{\sim}E_{tr}$ 是显然的同构（见 7.13.18 的注 (i)）。映射 $i^{\sim}=(i^{\sim}_+,i^{\sim}_-):F\widehat{\otimes}F^{\circ}(\infty\Delta)\to E_+\oplus E_-$ 把「核」$k$ 送到算子 $i^{\sim}_-(k)$ 与 $i^{\sim}_+(k)$：前者等于 $f\mapsto-\operatorname{Res}_{t_{2}=0}(k(t,t_{2}),f(t_{2}))dt_{2}$，后者等于 $f\mapsto(\operatorname{Res}_{t_{2}=t}+\operatorname{Res}_{t_{2}=0})(k(t,t_{2}),f(t_{2}))dt_{2}$。这里 $f\in F$，且 $(k(t,t_{2}),f(t_{2}))dt_{2}\in F((t_{2}))dt_{2}$。算子 $i^{\sim}_{\pm}(k)$ 属于 $E_{\pm}$ 一事留给读者验证\footnote{对 $i^{\sim}_-(k)$ 这是显然的。要验证 $i^{\sim}_+(k)\in E_+$，可用 Parshin 留数公式（[Pa76]，§1，命题 7）作用于 2-形式 $(k(t_{1},t_{2}),g(t_{1})f(t_{2}))dt_{1}\wedge dt_{2}$，其中 $g$ 属于 $F^{*}$ 中充分小的 c-格。}。因 $i^{\sim}$ 保持迹泛函，它给出所要的 $D^{Lie}$-模同构 $D^{\flat'}\xrightarrow{\sim}D^{\flat}$。\hfill$\Box$

\emph{注.} 设 $D_{i}\subset D$ 是次数 $\le i$ 的微分算子构成的空间。扩张 $D^{\flat}_{i}$ 带由嵌入\SourcePage{349}%
$D^{\flat}_{i}\subset\overline{\mathrm{Cl}}_{F}$ 引出的自然拓扑。这是 Tate 拓扑；$D_{i}$ 上的商拓扑与其作为有限维 $K$-向量空间的自然拓扑一致。

\bdp{7.13.21} \emph{例.} 置 $\mathcal{E}:=\operatorname{End}_{K}F=D_{0}\subset D$，于是得到 $\mathcal{E}^{Lie}$ 的中心扩张 $\mathcal{E}^{\flat}$。设 $\mathcal{L}\subset D^{Lie}$ 是 $\mathcal{E}$ 的正规化子；它以伴随的方式作用在 $\mathcal{E}^{\flat}$ 上。我们将把扩张 $\mathcal{E}^{\flat}$ 描述为 $\mathcal{L}$-模\footnote{因 $\mathcal{E}\subset\mathcal{L}$，这里特别描述了 $\mathcal{E}$ 的伴随作用，它正是 $\mathcal{E}^{\flat}$ 上的李括号。}。

易见 $\mathcal{L}$ 等于阶 $\le1$、其象征属于 $\operatorname{Der}_{K}\cdot\operatorname{id}_{F}$ 的微分算子李代数。换言之，$\mathcal{L}$ 由对 $(\tau,\tau^{\sim})$ 组成，其中 $\tau\in\operatorname{Der}K$、$\tau^{\sim}$ 是 $\tau$ 在 $F$ 上的作用；即 $\mathcal{L}$ 是 $(K,F)$ 的无穷小对称李代数。

同上，置 $\mathcal{E}^{\circ}:=\mathcal{E}\otimes_{K}\omega_{K}$。我们用配对 $\langle,\rangle:\mathcal{E}^{\circ}\times\mathcal{E}\to\mathbb{C}$、$\langle a,b\rangle:=\operatorname{Res}\operatorname{tr}_{K}(ab)$ 把 $\mathcal{E}^{\circ}$ 与 Tate 向量空间对偶 $\mathcal{E}^{*}$ 等同。$\mathcal{L}$ 在 $\mathcal{E}^{\circ}$ 上的伴随作用是 $(\tau,\tau^{\sim})(e\otimes\nu)=[\tau^{\sim},e]\otimes\nu+e\otimes\operatorname{Lie}_{\tau}\nu$。设 $\omega_{K}^{\otimes1/2}$ 是 $\operatorname{Spec}K$ 上的半形式层。它带 $\mathcal{L}$-作用（$(\tau,\tau^{\sim})$ 通过 $\operatorname{Lie}_{\tau}$ 作用），故 $\mathcal{L}$ 作用在 $\otimes\omega_{K}^{\otimes1/2}$ 上。考虑 $F\otimes\omega_{K}^{\otimes1/2}$ 上联络的集合 $\operatorname{Conn}(F\otimes\omega_{K}^{\otimes1/2})$\footnote{它不依赖于 $\omega_{K}^{\otimes1/2}$ 的选取。}。因 $\operatorname{End}_{K}F=\operatorname{End}_{K}(F\otimes\omega_{K}^{\otimes1/2})$，我们的 $\operatorname{Conn}(F\otimes\omega_{K}^{\otimes1/2})$ 是 $\mathcal{E}^{\circ}$-torsor（主齐性空间）；$\mathcal{L}$ 以显然的方式作用于其上。
% [Tutor Note: 原书 p.349 此处印作 "so $L$ acts on $\otimes\omega_{K}^{\otimes1/2}$"，$\omega$ 前多印一个 $\otimes$，底本如此，保留原样。]

\bdp{7.13.22} \emph{引理.} 存在唯一的 $\mathcal{L}$-且 $\mathcal{E}^{\circ}$-不变配对
\[
\langle,\rangle:\operatorname{Conn}(F\otimes\omega_{K}^{\otimes1/2})\times\mathcal{E}^{\flat}\to\mathbb{C}
\]
使对任意 $\nabla\in\operatorname{Conn}(F\otimes\omega_{K}^{\otimes1/2})$ 有 $\langle\nabla,1_{\mathcal{E}^{\flat}}\rangle=1$。

\emph{注.} (i) 元素 $\lambda\in\mathcal{E}^{\circ}$ 按公式 $\nabla\mapsto\nabla+\lambda$ 与 $e^{\flat}\mapsto e^{\flat}+\langle\lambda,e\rangle$ 分别作用在 $\operatorname{Conn}(F\otimes\omega_{K}^{\otimes1/2})$ 与 $\mathcal{E}^{\flat}$ 上（这里 $e:=e^{\flat}\bmod\mathbb{C}_{\mathcal{E}^{\circ}}=\mathcal{E}$）。故 $\langle,\rangle$ 的 $\mathcal{E}^{\circ}$-不变性即指 $\langle\nabla+\lambda,e^{\flat}\rangle=\langle\nabla,e^{\flat}\rangle-\langle\lambda,e\rangle$。

(ii) 显然 $\langle,\rangle$ 把 $\mathcal{E}^{\flat}$ 与 $\operatorname{Conn}(F\otimes\omega_{K}^{\otimes1/2})$ 上连续仿射泛函构成的 $\mathcal{L}$-模等同。这就是所承诺的 $\mathcal{E}^{\flat}$ 的描述。

\SourcePage{350}

\emph{证明.} $\langle,\rangle$ 的唯一性由下述事实推出：$\operatorname{Conn}(F\otimes\omega_{K}^{\otimes1/2})$ 无 $\mathcal{L}$-不变元素。

为定义 $\langle\nabla,e^{\flat}\rangle$，我们选取 $F$ 上的联络 $\nabla_{F}$ 与 $\omega_{K}$ 上的联络 $\nabla_{\omega}$，使得 $\nabla=\nabla_{F}+\frac{1}{2}\nabla_{\omega}$。

a. 联络 $\nabla_{F}$ 把 $F\otimes K$ 与 $K\otimes F$ 在对角线形式邻域上的限制等同，即它给出 $K\widehat{\otimes}K$-模的同构 $\epsilon(\nabla_{F}):F\widehat{\otimes}K\xrightarrow{\sim}K\widehat{\otimes}F$。令 $\epsilon(\nabla_{F}):F\widehat{\otimes}F^{\circ}\to K\widehat{\otimes}\omega_{K}$ 为 $\epsilon(\nabla_{F})\otimes\operatorname{id}_{F^{\circ}}$ 与显然态射 $K\widehat{\otimes}(F\otimes F^{\circ})\to K\widehat{\otimes}\omega_{K}$（由配对 $F\otimes F^{\circ}\to\omega_{K}$ 定义）的复合。用对角线方程局部化 $\epsilon(\nabla_{F})$，便得到态射 $F\widehat{\otimes}F^{\circ}(\infty\Delta)\to K\widehat{\otimes}\omega_{K}(\infty\Delta)$。把它作用于 $e^{\flat}$，得到 1-形式 $\epsilon(\nabla_{F},e^{\flat})\in K\widehat{\otimes}\omega_{K}(\Delta)$，它在那些满足 $\operatorname{Res}_{0}\phi(t,t)dt=0$ 的形式 $\phi(t_{1},t_{2})dt_{2}\in K\widehat{\otimes}\omega_{K}$ 构成的子空间模下良定义。注意对 $\lambda\in\mathcal{E}^{\circ}$ 有 $\epsilon(\nabla_{F}+\lambda,e^{\flat})=\epsilon(\nabla_{F},e^{\flat})-\operatorname{tr}_{K}(\lambda\cdot e)$（这里 $\operatorname{tr}_{K}(\lambda\cdot e)\in\omega_{K}=K\widehat{\otimes}\omega_{K}/(t_{1}-t_{2})K\widehat{\otimes}\omega_{K}$）。

b. 设 $\nu\in\omega_{K}\widehat{\otimes}K(\Delta)$ 是在对角线上留数为 1 的形式（即 $\nu$ 模 $\omega_{K}\widehat{\otimes}K$ 等于 $\frac{dt_{1}}{t_{1}-t_{2}}$）。设 $\psi(\nabla_{\omega})$ 是类似的形式，满足 $\psi(\nabla_{\omega})^{\otimes2}=-\nabla^{(1)}_{\omega}\nu$\footnote{这里 $\nabla^{(1)}_{\omega}$ 是沿第一个变量的协变导数。}。注意 $\psi(\nabla_{\omega})$ 在模 $(t_{1}-t_{2})\omega_{K}\widehat{\otimes}K$ 的意义下良定义。对 $l\in\omega_{K}$ 有 $\psi(\nabla_{\omega}+l)=\psi(\nabla_{\omega})-l$（这里把 $l$ 看作 $\omega_{K}\widehat{\otimes}K/(t_{1}-t_{2})\omega_{K}\widehat{\otimes}K$ 的元素）。

c. 考虑 2-形式 $\epsilon(\nabla_{F},e^{\flat})\wedge\nu$。置
\[
\langle\nabla,e^{\flat}\rangle:=\operatorname{Res}_{0}\operatorname{Res}_{\Delta}(\epsilon_{\nabla}(e^{\flat})\wedge\nu)
\]
则 $\langle\nabla,e^{\flat}\rangle$ 良定义（即它不依赖辅助选取），且 $\langle,\rangle$ 是 $\mathcal{E}^{\circ}$-不变的。因所有构造都是自然的，它也 $\mathcal{L}$-不变。\hfill$\Box$
% [Tutor Note: 原书 p.350 此处印作 "Since all the constructions where natural"，"where" 系 "were" 之误，底本如此，保留原样。]

\emph{注.} (i) 设 $e_{\alpha}$ 是 $F$ 的 $F$-基，$e'_{\alpha}$ 是 $F'$ 的对偶基，$\nabla$ 是使 $e'_{\alpha}\cdot(dt)^{-1/2}$ 为水平截面的联络。以 $(e_{\alpha}\cdot e'_{\beta})^{\flat}\in\mathcal{E}^{\flat}$ 记 $e_{\alpha}\otimes e'_{\beta}\frac{dt_{2}}{t_{2}-t_{1}}$ 的像。则 $\langle\nabla,(e_{\alpha}\cdot e_{\beta})^{\flat}\rangle=\delta_{\alpha,\beta}$。
% [Tutor Note: 原书 p.350 末式印作 "$(e_\alpha\cdot e_\beta)^{\flat}$"，第二处 $e_\beta$ 脱漏撇号（上文定义为 $(e_\alpha\cdot e'_\beta)^{\flat}$）；又上句印作 "an F-basis of F"，"F" 疑为 "K" 之误。底本如此，保留原样。]

\SourcePage{351}

(ii) 上述引理是局部 Riemann--Roch 定理的特例；参看 [BS] 的附录。

\bdp{7.13.23} 现在设 $\mathfrak{n}$ 是 Tate 意义下的李代数，即带连续李括号 $[\,,\,]$ 的 Tate 向量空间。下述引理或可使读者更安心。

\emph{引理.} $\mathfrak{n}$ 有由 $\mathfrak{n}$ 的李子代数构成的 $0$ 的邻域基。

\emph{证明.} 取任意 c-格 $P\subset\mathfrak{n}$。我们要找开李子代数 $\mathfrak{k}\subset P$。注意
\begin{equation}\tag{375}
\mathfrak{n}_{P}:=\{\alpha\in\mathfrak{n}:[\alpha,P]\subset P\}
\end{equation}
是开李子代数。置 $\mathfrak{k}:=P\cap\mathfrak{n}_{P}$。\hfill$\Box$

\bdp{7.13.24} 我们对 $F=\mathfrak{n}$ 使用 7.13.17 的记号。于是得到 Clifford 分次拓扑代数 $\overline{\mathrm{Cl}}^{\cdot}=\overline{\mathrm{Cl}}^{\cdot}_{\mathfrak{n}}$、相应的经典 Clifford 分次拓扑代数 $\overline{\mathcal{C}l}^{\cdot}=\operatorname{gr}\overline{\mathrm{Cl}}^{\cdot}$（它是泊松分次拓扑代数）、Tate 向量空间 $\mathfrak{n}$ 的自同态李代数 $E^{Lie}$ 的中心扩张 $E^{\flat}$，以及嵌入 $E^{\flat}\hookrightarrow\overline{\mathrm{Cl}}^{0}$。伴随作用给出态射 $\mathfrak{n}\to E^{Lie}$；以 $\mathfrak{n}^{\flat}$ 记扩张 $E^{\flat}$ 到 $\mathfrak{n}$ 的拉回。于是 $\mathfrak{n}^{\flat}$ 是 $\mathfrak{n}$ 被 $\mathbb{C}$ 的中心扩张。我们把 $\mathfrak{n}^{\flat}$ 赋以最弱的拓扑，使投影 $\mathfrak{n}^{\flat}\to\mathfrak{n}$ 与态射 $\mathfrak{n}^{\flat}\to\overline{\mathrm{Cl}}^{0}$ 都连续。则 $\mathfrak{n}^{\flat}$ 是 Tate 空间，且映射 $\mathfrak{n}^{\flat}/\mathbb{C}\to\mathfrak{n}$ 是同胚\footnote{确实，扩张 $\mathfrak{n}^{\flat}$ 在形如 (375) 的任意子代数上有典范连续分裂（其像由零化 $\lambda_{P}$ 的算子组成）。}。

\bdp{7.13.25} 现在我们准备把 BRST 构造搬到无穷维情形。先从「经典」版本开始。设 $\mathcal{R}$ 是拓扑泊松代数。我们假设 $\mathcal{R}$ 完备、分离，且是拓扑。
% [Tutor Note: 原书 p.351 末句作 "We assume that R is complete and separated and topology."，句子残缺，"and topology" 前疑有脱漏，底本如此，保留原样。]

\SourcePage{352}

\bdp{7.13.26} 以 $\mathcal{M}(\mathfrak{g})^{\flat}$ 记离散 $\mathfrak{g}^{\flat}$-模 $V$ 的范畴，其中 $1\in\mathbb{C}\subset\mathfrak{g}^{\flat}$ 以 $-\operatorname{id}_{V}$ 作用。对这样的 $V$，$\mathfrak{g}^{\flat}$ 在 $C^{\cdot}$ 与 $V$ 上的作用给出 $C^{\cdot}\otimes V$ 上的 $\mathfrak{g}$-模结构。它显然也是 $\mathrm{Cl}_{\mathfrak{g}}$-模，且 $\mathfrak{g}$-作用与 Clifford 作用相容。对 $\alpha\in\mathfrak{g}$，我们把其在 $C^{\cdot}\otimes V$ 上的作用记作 $\operatorname{Lie}_{\alpha}$，把 Clifford 算子 $C^{\cdot}\otimes V\to C^{\cdot-1}\otimes V$ 记作 $i_{\alpha}$。

宜把作用在 $C^{\cdot}\otimes V$ 上的算子改写如下（参看 7.7.5）。设 $\Omega_{\mathfrak{g}}$ 是 $\mathfrak{g}$ 的连续李代数上链构成的 DG 代数。相应的平面分次代数 $\Omega^{\cdot}_{\mathfrak{g}}$ 是 $\mathfrak{g}^{*}$ 的完备外代数。我们把它与完备 Clifford 代数 $\overline{\mathrm{Cl}}_{\mathfrak{g}}$ 中由 $\mathfrak{g}^{*}\subset\mathrm{Cl}_{\mathfrak{g}}$ 生成的闭子代数等同，于是 $\Omega^{\cdot}_{\mathfrak{g}}$ 以 Clifford 算子作用在 $C^{\cdot}\otimes V$ 上。现在设 $\mathfrak{g}_{\Omega}$ 是如下定义的 DG 李代数。其非零分量只有 $\mathfrak{g}^{0}_{\Omega}=\mathfrak{g}^{-1}_{\Omega}=\mathfrak{g}$，微分 $\mathfrak{g}^{-1}_{\Omega}\to\mathfrak{g}^{0}_{\Omega}$ 是 $\operatorname{id}_{\mathfrak{g}}$，$\mathfrak{g}^{0}_{\Omega}$ 上的括号就是 $\mathfrak{g}$ 的括号。回忆 $\mathfrak{g}_{\Omega}$ 作用在 $\Omega_{\mathfrak{g}}$ 上（即 $\mathfrak{g}^{0}_{\Omega}$ 按余伴随方式作用，$\mathfrak{g}^{-1}_{\Omega}$ 按「常数」导子作用）。分次李代数 $\mathfrak{g}^{\cdot}_{\Omega}$ 通过算子 $\operatorname{Lie}_{\alpha}$ 与 $i_{\alpha}$ 作用在 $C^{\cdot}\otimes V$ 上。故 $C^{\cdot}\otimes V$ 是分次 $(\Omega^{\cdot}_{\mathfrak{g}},\mathfrak{g}^{\cdot}_{\Omega})$-模。
% [Tutor Note: 原书 p.352 此处印作 "The corresponding plane graded algebra"，"plane" 疑为 "plain" 之误，底本如此，保留原样。]

\bdp{7.13.27} \emph{命题.} 存在唯一线性映射 $d:C^{\cdot}\otimes V\to C^{\cdot+1}\otimes V$，使对任意 $\alpha\in\mathfrak{g}$ 有 $\operatorname{Lie}_{\alpha}=di_{\alpha}+i_{\alpha}d$。有 $d^{2}=0$，且 $C_{\mathfrak{g}}(V):=(C^{\cdot}\otimes V,d)$ 是 DG $(\Omega_{\mathfrak{g}},\mathfrak{g}_{\Omega})$-模。

\emph{证明.} 唯一性。两个这样的 $d$ 之差是与一切 $i_{\alpha}$ 交换的算子。易见此算子代数与由 $\mathfrak{g}^{-1}_{\Omega}$ 与 $\operatorname{End}V$ 生成的闭子代数一致。因它没有正次数的算子，故毕。

类似论证表明 $(\Omega_{\mathfrak{g}},\mathfrak{g}_{\Omega})$ 的作用与微分相容，且 $d^{2}=0$（先证 $[d,\operatorname{Lie}_{\alpha}]=0$，再证其余性质）。

存在性。我们显式写出 $d$。设 $e_{i}$（$i\in I$）是 $\mathfrak{g}$ 的拓扑基（见 4.2.13），$e^{*}_{i}$ 是 $\mathfrak{g}^{*}$ 的对偶基。对半无限（关于 $\mathfrak{g}$）\SourcePage{353}%
子集 $A\subset I$，以 $\lambda_{A}\subset C^{\cdot}$ 记与由 $e_{a}$（$a\in A$）生成的 c-格相应的齐次线 $\lambda^{C}$（见 7.13.16）。换言之，$\lambda_{A}$ 是被 Clifford 算子 $e_{a}$、$e^{*}_{b}$（$a\in A$，$b\in I\backslash A$）零化的向量构成的子空间。我们的 $C^{\cdot}$ 是诸 $\lambda_{A}$ 的直和。注意对上述 $a,b$ 有 $e^{*}_{a}(\lambda_{A})=\lambda_{A\backslash a}$、$e_{b}(\lambda_{A})=\lambda_{A\cup b}$。
% [Tutor Note: 原书 p.353 此处印作 "the homoguenous line"，"homoguenous" 系 "homogeneous" 之误，底本如此，保留原样。]

置 $V_{A}:=\lambda_{A}\otimes V$；则 $C^{\cdot}\otimes V$ 是诸 $V_{A}$ 的直和。对 $c\in I$ 置 $L_{c}:=\operatorname{Lie}_{e_{c}}$、$i_{c}:=i_{e_{c}}$；对半无限子集 $A,A'$，以 $L^{A,A'}_{c}$、$i^{A,A'}_{c}$ 记这些算子从 $V_{A}$ 到 $V_{A'}$ 的 $A,A'$-分量。

设 $A,B$ 是半无限子集，满足 $|A|-|B|=1$（这里 $|A|-|B|:=|A\backslash(A\cap B)|-|B\backslash(A\cap B)|$）。任取 $a=a_{A,B}\in A\backslash(A\cap B)$（此集合非空）。以 $d^{A,B}$ 记复合 $V_{A}\to V_{B\cup a}\to V_{B}$，其中第一个箭头是 $L^{A,B\cup a}_{a}$，第二个是 Clifford 算子 $e^{*}_{a}$。易见以 $d^{A,B}$ 为分量的算子 $d:C^{\cdot}\otimes V\to C^{\cdot+1}\otimes V$ 定义正确（用到下述事实：对任意 $v\in V$，只有有限多个 $a\in A$ 使 $L_{a}(\lambda_{A}\otimes v)$ 非零）。
% [Tutor Note: 原书 p.353 此处作 "for any $v\in V$ and there is only finitely many $a\in A$ such that..."，"and" 之后疑有脱漏，底本如此，译文略去该 and。]

还需说明我们的 $d$ 满足命题的条件，即对任意 $c\in I$ 有 $[d,i_{c}]=L_{c}$。此事可直接计算验证；要点是 $[L_{a},i_{b}]$ 关于 $a,b$ 的斜对称性。细节留给读者。\hfill$\Box$

\bdp{7.13.28} 若 $V$ 是 $\mathcal{M}(\mathfrak{g})^{\flat}$ 中的复形，则以 $C_{\mathfrak{g}}(V)$ 记双复形 $C(V^{\cdot})$ 的全复形。这是离散 DG $(\Omega_{\mathfrak{g}},\mathfrak{g}_{\Omega})$-模（简称 $(\Omega_{\mathfrak{g}},\mathfrak{g}_{\Omega})$-复形）。函子 $C_{\mathfrak{g}}$ 是 DG 范畴 $C(\mathfrak{g})^{\flat}$（即 $\mathcal{M}(\mathfrak{g})^{\flat}$ 中的复形范畴，我们称之为 $\mathfrak{g}^{\flat}$-复形）与 $(\Omega_{\mathfrak{g}},\mathfrak{g}_{\Omega})$-复形构成的 DG 范畴 $C(\Omega_{\mathfrak{g}},\mathfrak{g}_{\Omega})$ 之间的等价。逆函子把 $F\in C(\Omega_{\mathfrak{g}},\mathfrak{g}_{\Omega})$ 送到复形 $\operatorname{Hom}_{\mathrm{Cl}_{\mathfrak{g}}}(C^{\cdot},F)$。

\bdp{7.13.29} 设 $\mathfrak{k}\subset\mathfrak{g}$ 是开有界李子代数。对 $a\ge0$，以 $C^{\cdot}_{a}\subset C^{\cdot}$ 记被来自 $\mathfrak{k}^{\perp}\subset\mathfrak{g}^{*}$ 的任意 $a+1$ 个 Clifford 算子的乘积零化的元素构成的子空间。则 $0=C^{\cdot}_{-1}\subset C^{\cdot}_{0}\subset C^{\cdot}_{1}\subset\ldots$ 是 $C^{\cdot}=\cup C^{\cdot}_{a}$ 上的递增滤链。$\mathfrak{g}^{*}$ 中任意 Clifford 算子 $\nu$ 保持\SourcePage{354}%
我们的滤链；若 $\nu$ 属于 $\mathfrak{k}^{\perp}$，则它把 $C^{\cdot}_{a}$ 送到 $C^{\cdot+1}_{a-1}$。来自 $\mathfrak{g}$ 的任意 Clifford 算子把 $C^{\cdot}_{a}$ 送到 $C^{\cdot-1}_{a+1}$；若它属于 $\mathfrak{k}$，则保持滤链。于是 $\operatorname{gr}_{*}C^{\cdot}$ 是向量空间 $(\mathfrak{g}/\mathfrak{k})\oplus(\mathfrak{g}/\mathfrak{k})^{*}\oplus\mathfrak{k}\oplus\mathfrak{k}^{*}$（赋以标准「双曲」形式）的 Clifford 代数 $\mathrm{Cl}_{\mathfrak{g},\mathfrak{k}}$ 上的模。

这是不可约 $\mathrm{Cl}_{\mathfrak{g},\mathfrak{k}}$-模；且 $C^{\cdot}_{0}$ 是子代数 $\mathrm{Cl}_{\mathfrak{k}}\subset\mathrm{Cl}_{\mathfrak{g},\mathfrak{k}}$ 上的不可约模。齐次线 $\lambda_{\mathfrak{k}}=\lambda^{(C)}_{\mathfrak{k}}$（见 7.13.16）位于 $C^{\cdot}_{0}$ 中，且 $\operatorname{gr}_{*}C^{\cdot}$ 是由此线生成的、作为子代数 $\Lambda(\mathfrak{g}/\mathfrak{k})\otimes\Lambda\mathfrak{k}^{*}\subset\mathrm{Cl}_{\mathfrak{g},\mathfrak{k}}$ 上的自由模。若 $\lambda_{\mathfrak{k}}\subset C^{0}$（必要时可平移 $\cdot$-滤链来假定这一点），则 $\operatorname{gr}_{a}C^{b}=\Lambda^{a}(\mathfrak{g}/\mathfrak{k})\otimes\Lambda^{b+a}\mathfrak{k}^{*}\otimes\lambda_{\mathfrak{k}}$。
% [Tutor Note: 原书 p.354 此处作 "(we may assume this shifting the · filtration if necessary)"，语句不顺，"this shifting" 疑为 "this by shifting" 之误，又 "homoguenous" 系 "homogeneous" 之误；底本如此，保留原样。]

设 $\mathfrak{k}^{\flat}\subset\mathfrak{g}^{\flat}$ 是 $\mathfrak{k}$ 的原像。这是 $\mathfrak{k}$ 被 $\mathbb{C}$ 的中心扩张，它典范分裂：分裂 $\mathfrak{k}\to\mathfrak{k}^{\flat}$ 的像由那些零化 $\lambda_{\mathfrak{k}}$ 的元素组成（我们考虑 $\mathfrak{k}^{\flat}$ 在 $C^{\cdot}$ 上的李代数作用）。

对 $V\in C(\mathfrak{g})^{\flat}$，子空间 $C^{\cdot}_{a}\otimes V$ 是 $C_{\mathfrak{g}}(V)$ 的子复形；以 $C_{\mathfrak{g}}(V)_{a}$ 记之。我们得到 $C_{\mathfrak{g}}(V)$ 上的滤链，它被来自 $\mathfrak{g}^{*}$ 与 $\mathfrak{k}$ 的 Clifford 算子保持；相继商 $\operatorname{gr}_{a}C_{\mathfrak{g}}(V)$ 是 $(\Omega_{\mathfrak{k}},\mathfrak{k}_{\Omega})$-复形。对 $\mathfrak{k}$-复形 $P$，以 $C_{\mathfrak{k}}(P)$ 记 $\mathfrak{k}$ 的、系数在 $P$ 中的李代数上链构成的 Chevalley 复形；这是 $(\Omega_{\mathfrak{k}},\mathfrak{k}_{\Omega})$-复形。等同
\[
\operatorname{gr}_{a}C_{\mathfrak{g}}(V)^{\cdot}=\Lambda^{\cdot+a}\mathfrak{k}^{*}\otimes(V^{\cdot}\otimes\Lambda^{a}(\mathfrak{g}/\mathfrak{k})\otimes\lambda_{\mathfrak{k}})
\]
是 $(\Omega_{\mathfrak{k}},\mathfrak{k}_{\Omega})$-复形的同构
\begin{equation}\tag{376}
\operatorname{gr}_{a}C_{\mathfrak{g}}(V)\simeq C_{\mathfrak{k}}(V\otimes\Lambda^{a}(\mathfrak{g}/\mathfrak{k})\otimes\lambda_{\mathfrak{k}})[a]
\end{equation}

这里 $\mathfrak{k}$ 以伴随的方式作用在 $\Lambda^{a}(\mathfrak{g}/\mathfrak{k})$ 上。相应的谱序列收敛到 $H^{\cdot}C_{\mathfrak{g}}(V)$；其首项是 $E^{p,q}_{1}=H^{p+q}\operatorname{gr}_{-p}C_{\mathfrak{g}}(V)=H^{q}(\mathfrak{k},\Lambda^{-p}(\mathfrak{g}/\mathfrak{k})\otimes V\otimes\lambda_{\mathfrak{k}})$。

\bdp{7.13.30} \emph{注.} 假设有 $\mathfrak{k}^{\flat}$-子复形 $T\subset V$，使 $V$ 由 $T$ 诱导出，即 $V=U(\mathfrak{g}^{\flat})\otimes_{U(\mathfrak{k}^{\flat})}T$。则嵌入的复合 $C_{\mathfrak{k}}(T\otimes\lambda_{\mathfrak{k}})\subset C_{\mathfrak{g}}(V)_{0}\subset C_{\mathfrak{g}}(V)$ 是拟同构。

\SourcePage{355}

\bdsec{7.14. 无穷维情形下的局部化函子} 现在我们可以解释 7.1.1 中「Hecke 模式」的部分 (c)、(d) 在当前的无穷维情形下是什么意思。

\bdp{7.14.1} 设 $G,K$ 如 7.11.17，$G'$ 是 $G$ 被 $\mathbb{G}_{m}$ 的中心扩张，并带分裂 $K\to G'$（参看 7.8.1）。则 $\mathfrak{g},\mathfrak{g}'$ 是 Tate 意义下的李代数，而 $\mathfrak{k}=\operatorname{Lie}K$ 是 $\mathfrak{g},\mathfrak{g}'$ 的开有界李子代数。7.8.1 中所有范畴在当前情形下都有显然的意义。

按 7.8.2 的方式在范畴 $D(\mathfrak{g},K)'$ 上定义 Hecke 作用。现在线丛 $\mathcal{L}_{G}$ 是 $G$ 上的 $\mathcal{O}^{p}$-模，而 $\mathcal{V}_{G}$ 是左 $\mathcal{D}^{p}$-模构成的复形（见 7.11.3）。7.8.2 的所有构造逐字搬到我们的情形，7.8.4--7.8.5 亦然（在 7.8.4 中，$U'$ 照例应取完备化扭包络代数）。

\bdp{7.14.2} 为定义局部化函子 $L\Delta$，我们需要一些准备。设 $Y$ 是概形，$F$ 是 Tate 向量空间。$Y$ 上的 $\mathrm{Cl}_{F}$-模是 $Y$ 上的 $\mathbb{Z}$-分次 $\mathcal{O}$-模 $C^{\cdot}$，并带分次 Clifford 代数 $\mathrm{Cl}^{\cdot}_{F}$ 的连续作用（见 7.13.16）。对任意 c-格 $P\subset F$，以 $\lambda_{P}(C^{\cdot})$ 记 $C^{\cdot}$ 的如下分次 $\mathcal{O}$-子模：它由被来自 $P\subset F$ 与 $P^{\perp}\subset F^{*}$ 的 Clifford 算子零化的局部截面组成。函子 $\lambda_{P}:\mathcal{C}(Y)\to\{Y$ 上分次 $\mathcal{O}$-模的范畴$\}$ 是范畴等价\footnote{逆函子是与适当的不可约分次 Clifford 模作张量积。}。对两个 c-格 $P_{1},P_{2}$，有典范同构
\begin{equation}\tag{377}
\lambda_{P_{1}}(C^{\cdot})\simeq\lambda_{(P_{1},P_{2})}\otimes\lambda_{P_{2}}(C^{\cdot})
\end{equation}
它满足显然的传递性（见 7.13.16）。对 c-格的 $Y$-族亦然（见 loc. cit.）。

\bdp{7.14.3} 现在假设处于 7.11.18 的情形。则 $Y$ 带典范的 $\mathrm{Cl}_{\mathfrak{g}}$-模 $C^{\cdot}_{Y}$，其定义如下。设 $K\subset G$ 是合理的群子概形，$\mathfrak{k}:=\operatorname{Lie}K$。以 $\omega_{(K\backslash Y)}$ 记典范
% 尾部接 p.356（本单元在此断处停止）
